\begin{frame}{Motivation}

\begin{block}{\faBuildingColumns~Academia}
Aims to prepare students for professional practice, but offers limited guidance on how to operationalize this in course design and assessment
\end{block}

\begin{block}{\faIndustry~Industry}
Expects credible, reproducible results supported by clear reasoning and usable for decision-making
\end{block}

\begin{alertblock}{\faScaleBalanced~Underlying mismatch}
What counts as ``good work'' differs between academic and professional contexts
\end{alertblock}

\end{frame}
\note{
This work starts from a familiar gap.  

Universities aim to prepare students for professional practice --  
but translating that into course design is not straightforward.  

At the same time, industry expectations are quite concrete.  
Work needs to be credible, reproducible, and usable.  

And these perspectives do not always align.  

So the question is not whether we value the same things --  
but how those values are actually operationalized in courses.
}

%------------------------------------------------
\begin{frame}{What prompted the redesign}

\begin{alertblock}{Observed in an industry-facing course}
\begin{itemize}
    \item[\faIndustry] Industry partners valued collaboration -- but would not use the results
    \item[\faUsers] Students reported a very heavy and uneven workload
\end{itemize}
\end{alertblock}

\begin{exampleblock}{Core question}
How do we design courses so that student work is \emph{authentic in practice},\\
while still supporting \emph{meaningful learning}?
\end{exampleblock}

\end{frame}
\note{
This is not just a theoretical issue.  

In an earlier version of the course, we worked with an industry partner  
and asked how useful the student results were.  

They appreciated the collaboration --  
but they would not actually use the results.  

At the same time, students experienced the course as quite heavy  
and time-consuming.  

So the course was not fully working for either side.  

That leads to the core question:  

how do we design assessment so that the work is \emph{authentic in practice},  
while still supporting \emph{meaningful learning}?
}

